\subsection{오케스트레이션 구조의 효율성 및 품질 트레이드-오프 분석}

본 연구의 실험 결과는 오케스트레이션 제어 메커니즘의 설계 방식이 시스템의 자원 소모와 생성 품질 사이에서 밀접한 트레이드-오프(Trade-off) 관계를 형성함을 보여준다. 고정 그래프 기반의 Code 구조는 정형 보고서 생성 작업에서 평균 실행 시간 158.61초, 토큰 소모량 21,440개를 기록하였다. 이러한 수치는 결정론적인 제어 흐름을 가짐에도 불구하고, 실제 정형 보고서를 생성하기 위한 개별 태스크의 복잡도로 인해 상당한 수준의 자원 소모가 발생함을 시사한다.

이론적으로 LLM 기반의 동적 라우팅이나 Hybrid 구조는 제어 단계에서 LLM이 현재 상태를 분석하고 다음 에이전트를 결정하는 추론 과정이 추가되므로, Code 구조보다 실행 시간과 토큰 비용이 더욱 증가할 가능성이 높다. 그러나 품질 관점에서 볼 때, 고정된 파이프라인은 예외적인 입력 데이터나 복잡한 문맥적 수정 요구사항에 유연하게 대응하지 못하는 한계가 존재한다. 따라서 시스템 설계자는 단일한 효율성 지표에 의존하기보다, 작업의 목적에 따라 비용과 품질 간의 균형을 고려한 구조 선택이 필요하다. 예를 들어, 정해진 양식의 준수가 최우선인 산업 보고서의 경우 Code 구조가 적절한 선택지가 될 수 있으나, 창의적 구성과 반복적인 정교화가 필요한 고부가가치 콘텐츠 생성에서는 자원 소모 증가를 감수하더라도 Hybrid 또는 LLM 구조를 채택하는 것이 타당하다.

\subsection{실무적 설계를 위한 오케스트레이션 선택 가이드라인}

실험을 통해 도출된 정량적 지표를 바탕으로, 실제 멀티 에이전트 시스템을 구축하려는 설계자를 위한 의사결정 매트릭스를 제안한다. 오케스트레이션 구조의 선택 기준은 크게 작업의 정형성(Formality)과 요구되는 복잡도(Complexity)라는 두 가지 축으로 정의될 수 있다.

표~\ref{tab:guideline}에 제시된 가이드라인과 같이, 작업의 성격에 따라 다음과 같은 설계 전략을 권고한다. 첫째, 입력 데이터의 구조가 명확하고 출력 양식이 고정된 '정형 보고서' 작업의 경우, Code 구조를 통해 제어 오버헤드를 방지하고 프로세스를 정형화하는 것이 효율적이다. 둘째, 작업의 단계가 가변적이며 에이전트 간의 상호 피드백 루프가 빈번하게 발생해야 하는 '창의적 글쓰기'나 '심층 분석' 작업의 경우, LLM 기반의 동적 라우팅이 필수적이다. 셋째, 전체적인 흐름은 정해져 있으나 세부 단계에서 품질 검수 및 재시도(Retry) 로직이 정밀하게 제어되어야 하는 경우, 계층적 제어를 수행하는 Hybrid 구조가 적절한 절충안이 된다.

이러한 가이드라인은 작업의 성격에 따라 자원 소모와 품질 간의 균형을 고려한 설계 전략을 제시하며, 설계자가 가용한 컴퓨팅 자원과 허용 가능한 지연 시간, 그리고 최종 결과물의 품질 임계치를 종합적으로 판단하여 구조를 선택하도록 돕는다.

\subsection{연구의 한계 및 향후 발전 방향}

본 연구는 오케스트레이션 구조의 성능 차이를 정량적으로 분석하였으나, 몇 가지 한계점을 갖는다. 첫째, 품질 평가 지표로 사용된 LLM-as-Judge 방식은 평가 모델의 내재적 편향(Bias)에서 완전히 자유로울 수 없으며, 정량적 점수 외에 인간 전문가의 정성적 평가를 통한 교차 검증이 추가적으로 필요하다. 둘째, 실험에 사용된 LLM 모델의 종류와 버전이 고정되어 있어, 모델의 추론 능력 변화에 따른 오케스트레이션 효율성의 변동성을 완전히 포괄하지 못했다.

향후 연구에서는 오케스트레이션 구조를 정적으로 선택하는 것을 넘어, 작업의 난이도나 초기 생성물의 품질에 따라 실시간으로 제어 구조를 변경하는 '적응형 오케스트레이션(Adaptive Orchestration)' 메커니즘을 탐구할 계획이다. 또한, 토큰 소모량을 더욱 획기적으로 줄이기 위해 제어 전용 소형 언어 모델(sLLM)을 오케스트레이터로 배치하는 구조적 최적화 방안에 대해서도 심도 있게 분석하고자 한다.